\section{Act II --- Does the objective matter? Design and pre-registration}
\label{sec:actII}

\begin{pending}{\textsc{pending} the locked Paper C GPU run --- design and hypotheses only, no numbers invented}
This section states the \emph{question, objective menu, reward design, comparability contract,
pre-registered hypotheses, and analysis plan} for the objective axis in full. It reports \textbf{no
results}: the DPO/GRPO (and extension KTO/ORPO) training runs are executed only after this design is
bound into the Paper C lock, and their per-row scores are not yet committed. Every prediction below is
a \emph{pre-registration}, fixed in advance (\code{docs/paper-c-prereg.md}, whose SHA-256 is recorded
in the lock \emph{before} the first final run so the outcome cannot be reverse-fit). Like Acts I and
III, the evidence tier is \emph{retrospective, estimation-only, fixed-panel --- not confirmatory}. The
results table will be \code{\textbackslash input} here as \code{tab:paper-c} once the locked run
completes (\Cref{sec:actII-results}).
\end{pending}

\subsection{The question}
\label{sec:actII-q}

Act I (\Cref{sec:actI}) held one thing constant that practitioners routinely vary without a second
thought: the \emph{training objective}. It fine-tuned every guard with supervised fine-tuning (SFT) ---
cross-entropy on the correct verdict token --- and found the specialization signature (represented-source
ranking up to a $\approx\!0.98$ ceiling for every checkpoint, transfer flat-to-down; \Cref{tab:primary},
\Cref{fig:plane}). The natural next question is whether that trade-off is a property of \emph{fine-tuning
the guard at all}, or a property of the \emph{particular objective} SFT. If a different objective moved
the model to the same represented ceiling while giving up less transfer, ``how you tune'' would matter as
much as ``whether you tune,'' and the specialization tax could be partly a choice rather than a law.

\begin{background}{What is a ``fine-tuning objective,'' and why might it change the trade-off?}
All of these methods take the same base model and the same labeled prompts and end with a guard that
emits \code{safe}/\code{unsafe}. They differ only in the \emph{loss they minimize} --- the mathematical
statement of ``what counts as improvement.'' \textbf{SFT} simply maximizes the probability of the correct
verdict token. \textbf{Preference} methods (DPO, KTO, ORPO) instead push the correct verdict \emph{above}
the wrong one by a margin, some while tethering the model to its starting point with a Kullback--Leibler
(KL) ``leash.'' \textbf{Reinforcement-learning} methods (GRPO) let the model sample a verdict and reward
it when the sample is correct. The reason the objective \emph{could} change transfer: each loss moves the
model a different \emph{distance} from its base, and --- as the recent ``RL generalizes, SFT memorizes''
and ``RL's Razor'' analyses argue for multi-token tasks --- how far a model is dragged from its base is a
leading predictor of how much prior, off-distribution behavior it forgets. Our thesis for this axis is
that the same distance-from-base variable orders the specialization/transfer trade-off here too.
\end{background}

Because the represented-source AP saturates near $0.98$ for \emph{any} objective that can fit a
binary verdict (Act I already shows SFT does), the informative axis is \textbf{transfer} --- dataset-held-out
macro-AP --- and the mediating quantity is \textbf{how far each objective moves the tuned model from its
base}. The entire section is organized around that mediator.

\subsection{The objective menu, and why each arm earns its place}
\label{sec:actII-menu}

We select objectives \emph{mechanistically}, so that each arm anchors a distinct point on the
reference/KL (distance-from-base) axis rather than being chosen by taste. The menu splits three ways ---
pointwise, pairwise, and reward-based --- and each family is represented.

\begin{itemize}[leftmargin=1.4em,itemsep=1pt]
  \item \textbf{SFT (pointwise, no reference)} --- the field-standard guard recipe: Llama Guard,
  WildGuard, and ShieldGemma are all supervised instruction fine-tunes emitting a safe/unsafe verdict
  \citep{inan2023llamaguard,han2024wildguard,zeng2024shieldgemma}. It carries no KL leash, so it is the
  ``moves-most, reference-free'' endpoint and our reused Act I arm.
  \item \textbf{KTO (pointwise binary; \citealp{ethayarajh2024kto})} --- trains directly on a per-example
  \emph{desirable}/\emph{undesirable} label, which is \emph{exactly} the shape of our data, and carries
  built-in class-imbalance weights. It is KL/reference-anchored (a $\beta$ leash to the initial policy),
  so it is the ``most natural fit, anchored'' arm.
  \item \textbf{DPO (pairwise, reference-anchored; \citealp{rafailov2023dpo})} --- the canonical preference
  objective; we synthesize the required \emph{(chosen, rejected)} pair losslessly from the binary label
  (chosen $=$ correct verdict token, rejected $=$ wrong token) against a frozen base reference with knob
  $\beta$. Its bounded-loss sibling \textbf{IPO} \citep{azar2023ipo} is held as an optional robustness
  cross-check, included only if DPO shows the near-deterministic-pair overfitting the theory predicts.
  \item \textbf{ORPO (pairwise, reference-free; \citealp{hong2024orpo})} --- a single-stage odds-ratio
  penalty on the SFT loss that needs \emph{no} reference model. It removes the $\beta$ leash entirely, so
  it anchors the opposite (maximal-movement) end of the mechanism from the $\beta$-anchored arms and is
  scientifically load-bearing for H2.
  \item \textbf{GRPO (reward / RLVR; \citealp{shao2024deepseekmath})} --- the headline ``does RL
  generalize?'' test. Per prompt it samples a group of verdicts, and the advantage is the reward
  normalized within the group, with an optional KL-to-reference term. The label supplies the reward
  (\Cref{sec:actII-reward}).
\end{itemize}

\paragraph{What we deliberately drop, and why.} \textbf{SimPO}'s entire method is length-normalization,
which is \emph{moot at output length 1}: it collapses to a plain one-token margin loss nearly identical to
a margin DPO, so it adds no distinct axis point. \textbf{PPO} \citep{schulman2017ppo} is strictly
dominated by GRPO for a verifiable single-token reward and is heavier. \textbf{BCO} is a KTO sibling,
redundant with KTO for a first study. These choices trade breadth for a clean, mechanism-spanning menu.

\begin{table}[htbp]\centering\small\setlength{\tabcolsep}{5pt}
\caption{The objective menu for Act II. Each arm anchors a distinct point on the distance-from-base
(reference/KL) axis; the score at evaluation is \emph{identical} across every arm
($s=z_{\text{unsafe}}-z_{\text{safe}}$, one forward pass). The last column is a pre-registered
\emph{prediction} (H1, \Cref{sec:actII-hyp}), not a result. Locked minimal-viable core: SFT/DPO/GRPO;
KTO/ORPO (and a $\beta$ sweep) are the pre-specified extension.}
\label{tab:objective-menu}
\begin{tabular}{@{}llp{0.30\textwidth}lp{0.16\textwidth}@{}}
\toprule
Arm & Family & Signal built from the binary label & Ref./KL anchor & Predicted movement (H1) \\
\midrule
SFT  & pointwise    & cross-entropy on the correct verdict token & none            & most (tied) \\
KTO  & pointwise    & $\{$prompt, correct token, label$\}$ (optionally add wrong token) & yes ($\beta$) & least (anchored) \\
DPO  & pairwise     & pair: chosen $=$ correct token, rejected $=$ wrong token & yes ($\beta$, frozen base) & least (anchored) \\
ORPO & pairwise     & same pair; odds-ratio penalty on the SFT loss & \textbf{none}   & most (tied) \\
GRPO & reward (RLVR)& reward $=$ correctness of the sampled verdict vs.\ label & yes / optional & least (KL-close) \\
\bottomrule
\end{tabular}
\end{table}

\paragraph{Scope of the locked run vs.\ the full menu.} The pre-registered \emph{minimal-viable} run
that this section commits comprises \code{base} (reused), \code{sft} (reused Act I arm), \code{dpo}, and
\code{grpo} --- the three-way SFT/DPO/GRPO comparison named in the abstract and intro. \code{kto},
\code{orpo}, the $\beta$ sweep, and the IPO/RLOO cross-checks are the pre-specified Step-3 extension; they
require no new training campaign beyond adding their conditions, and they sharpen the mechanism (the
reference-free ORPO endpoint and the natural-fit KTO arm), but the core hypotheses are answerable with the
minimal set.

\subsection{The reward is the label: verifiable rewards, no learned reward model}
\label{sec:actII-reward}

\begin{background}{RLVR --- reinforcement learning with \emph{verifiable} rewards}
In general RLHF one trains a neural \emph{reward model} to imitate human preferences, then optimizes the
policy against it \citep{ouyang2022instructgpt}. That reward model is an approximation, and optimizing
against an approximation invites \emph{reward hacking} --- the policy finds inputs the reward model scores
highly but a human would not. When the task has a \emph{checkable} answer (a math result, or here a gold
safety label), the canonical modern recipe skips the neural reward model entirely and rewards
\emph{ground-truth correctness}. Fewer moving parts, no learned reward surface to exploit.
\end{background}

A safety guard has verifiable ground-truth labels, so \textbf{there is no reward model to design ---
the reward \emph{is} the gold label}. Introducing a learned reward model would be unnecessary and would
add exactly the reward-hacking surface that verifiable-reward recipes deliberately avoid. Crucially, this
keeps the \emph{evaluation identical across every arm}: whatever the training loss, the guard is always
scored by the same one-forward-pass quantity $s(x)=z_{\text{unsafe}}-z_{\text{safe}}$ and its two-way
softmax (\Cref{sec:setup}). No arm is scored by a sampled string; the arms are compared on the same
metric. ``RL for guards'' already exists in a \emph{different} sense --- e.g.\ DuoGuard uses RL for
adversarial \emph{data generation} while the guard stays a classifier \citep{deng2025duoguard} --- whereas
here the objective is applied to the guard's \emph{own verdict}.

Each objective's training signal is constructed deterministically from the same frozen label (the builder
is seeded and its \code{preference\_recipe\_sha256} is bound into the lock, so the preference/reward data is
reproducible): SFT uses the correct verdict token as the cross-entropy target; DPO/ORPO/IPO use the
(correct, wrong) verdict-token pair; KTO uses the pointwise $\{$completion, desirable/undesirable$\}$
record; GRPO uses the graded margin below.

\paragraph{The one genuine reward-design problem: GRPO on a single token.} A one-token, two-outcome action
space breaks GRPO's group-relative advantage. Once the guard is confident, all $G$ sampled verdicts agree,
the within-group standard deviation collapses to zero, the advantage vanishes, and that prompt contributes
\emph{no gradient} --- ``gradient starvation.'' Our design fixes this in three pre-registered ways.

\emph{(1) Graded margin instead of raw $0/1$ --- the single most important GRPO choice.} Rather than a
binary correctness reward, we reward the \emph{margin toward the correct side}, i.e.\ the calibrated
probability mass the guard places on the correct verdict:
\begin{equation}
r(x) \;=\; \sigma\!\big(s_y \cdot s(x)\big) \;=\; P_\theta(\text{correct verdict}\mid x),
\qquad
s_y=\begin{cases}+1 & y=\text{unsafe}\\[-1pt] -1 & y=\text{safe}\end{cases}
\label{eq:grpo-reward}
\end{equation}
where $s(x)=z_{\text{unsafe}}-z_{\text{safe}}$ and $\sigma$ is the logistic function. Because $r(x)$ is
continuous, it restores the within-group spread that a raw $0/1$ reward destroys, and it rewards the
\emph{exact} quantity the guard is evaluated on.

\emph{(2) Unnormalized (Dr.GRPO/RLOO-style) advantage.} With only two underlying reward values,
group-standard-deviation normalization is pure difficulty bias, so we drop it: for a group of $G$
sampled verdicts,
\begin{equation}
A_i \;=\; r_i - \tfrac{1}{G}\textstyle\sum_{j=1}^{G} r_j
\qquad\text{(rather than } A_i=(r_i-\mu)/\sigma\text{)}.
\label{eq:grpo-adv}
\end{equation}

\emph{(3) Dynamic sampling and class weighting.} Zero-variance (agreed-verdict) groups are filtered so
they add no noise; sampling temperature can be raised for diversity ($G>1$ is cheap since rollouts are one
token). The only live reward-hacking mode --- length hacking being impossible at one token --- is
majority-class collapse, guarded by a class/severity-weighted correctness reward if the frozen manifest is
skewed (the v2 audit asserts label balance, so any skew is small).

\begin{takeaway}{Why we predict this fix mostly recovers the \emph{same} signal}
The graded reward is honest, but it is telling: on a single verdict token there is no reasoning trace to
diversify, so \Cref{eq:grpo-reward} largely hands GRPO the same margin information SFT and KTO already
optimize. That is precisely why we \emph{pre-register} GRPO's transfer edge as likely small or null
(\Cref{sec:actII-hyp}) --- a bounded negative result, not a bug.
\end{takeaway}

\subsection{Held fixed for comparability}
\label{sec:actII-fixed}

The objective is the \emph{only} thing that varies. Everything that could otherwise confound a
cross-objective comparison is reused from Act I \emph{by hash}, so that a difference in transfer is
attributable to the loss and not to data, panel, or scoring drift:
\begin{itemize}[leftmargin=1.4em,itemsep=1pt]
  \item \textbf{Panel:} the identical four checkpoints (Qwen2.5-1.5B, SmolLM2-1.7B, SmolLM3-3B, Qwen3-4B).
  \item \textbf{Data:} the same frozen, decontaminated 1{,}200-row \code{train.jsonl} (400 rows each from
  ToxicChat / Prompt-Injections / Jailbreak-Classification, at a 200:200 safe:unsafe balance) and the same
  calibration / represented-ID / transfer / OR-Bench / HarmBench manifests, all consumed by SHA-256. No new
  data.
  \item \textbf{LoRA recipe \citep{hu2021lora}:} $r{=}32$, $\alpha{=}64$, dropout $0.05$, targets
  $\{$q,k,v,o,gate,up,down$\}$; max length $1024$; $300$ steps; lr $2\!\times\!10^{-4}$ cosine with $0.03$
  warmup; effective batch $4$.
  \item \textbf{Seeds:} 42--46 (five per arm), data-order seed 42 --- matching Paper A's $4\times5$ protocol
  so the bootstrap variance bars are directly comparable.
  \item \textbf{Guard head \& score:} the single next-token verdict scored by
  $s=z_{\text{unsafe}}-z_{\text{safe}}$, calibrated by the same development-split calibrator
  \citep{guo2017calibration}.
  \item \textbf{Provenance:} a fresh artifact root \code{artifacts/paper\_c\_objective\_v2/} with its own
  fail-closed lock that reuses Paper A's manifest/audit bindings by hash and records the objectives recipe,
  the \code{preference\_recipe\_sha256}, exact software versions, and this section's pre-registration
  SHA-256. Paper A's artifacts are never modified.
\end{itemize}
The reference-anchored arms (DPO/IPO/KTO, and GRPO-with-KL) hold a second frozen model; on a single 24\,GB
GPU this stays in budget via precomputed reference log-probabilities or a frozen LoRA-base reference, while
ORPO needs no second model. Verdict rollouts are one token, so GRPO's per-step cost is trivial and the
feasibility risk is convergence steps, not rollout length.

\subsection{Pre-registered hypotheses}
\label{sec:actII-hyp}

\begin{background}{Why pre-register?}
Fixing the hypotheses, metrics, and decision rules \emph{before} seeing the results prevents ``HARKing''
(Hypothesizing After the Results are Known) --- quietly retrofitting the story to whatever the data
happened to show. The predictions below are bound into the lock; any later change is a new, dated
revision, not a silent edit.
\end{background}

The mediator variable is \textbf{how far the objective moves the tuned model from its base}, measured on
the \emph{represented} rows as the forward KL (or, equivalently for our purposes, the shift in the
calibrated $P(\text{unsafe})$ distribution) between tuned and base:
\begin{equation}
M(\text{base}\!\to\!\text{tuned}) \;=\; \mathbb{E}_{x\,\in\,\text{represented}}
\Big[\,\mathrm{KL}\big(P_{\text{tuned}}(\cdot\mid x)\,\big\|\,P_{\text{base}}(\cdot\mid x)\big)\Big].
\label{eq:movement}
\end{equation}

\paragraph{H1 --- the objective orders the trade-off by movement from base.} Because every objective can
fit the verdict, represented-source AP rises to the same $\approx\!0.98$ ceiling for all arms (as SFT
already does; \Cref{tab:primary}); the arms therefore differ mainly on \emph{transfer}, and that
difference \emph{tracks} the movement mediator of \Cref{eq:movement}: objectives that move the model
\emph{more} lose \emph{more} transfer. The pre-registered movement order is
\[
\text{ORPO (reference-free)} \;\gtrsim\; \text{SFT} \;>\; \text{KTO} \approx \text{DPO/IPO ($\beta$-anchored)} \;>\; \text{GRPO (KL-close, low-signal)},
\]
with transfer loss following the same order. This direction is already \emph{directionally} visible in the
legacy single-seed table (SFT buys the most in-distribution gain and risks transfer; GRPO barely moves off
base; DPO sits between); the clean rerun's job is to confirm it under the corrected macro-AP with seed
variance.

\paragraph{The GRPO single-token null (the headline pre-registration).} The ``RL generalizes'' results
that would predict GRPO an \emph{edge} over SFT were established for \emph{multi-token generative reasoning}
tasks, where outcome-reward exploration over long traces is the source of the gain. Our action space is
\emph{one token, two outcomes}, which removes those exploration/reasoning degrees of freedom. We therefore
predict \textbf{GRPO's transfer edge over SFT will be small or null} --- and we state this up front so that
a bounded negative result is the \emph{expected, reportable} outcome rather than a disappointment. Reported
honestly, it \emph{bounds where the RL-generalizes claim applies}. This framing is consistent with the
recent finding that constrained objectives can help general LLM safety at scale
\citep{vennemeyer2026objective}, which validates the ``objective matters'' framing without predicting a win
for RL on a single-token classifier.

\paragraph{H2 --- objective $\times$ composition interaction (the genuinely new question).} The
output-space composition of \Cref{sec:actIII} (\Cref{eq:comp}, averaging the base and tuned calibrated
probabilities) is an output-space analog of weight-space WiSE-FT / model-soup robustness recovery
\citep{wortsman2022wiseft,wortsman2022modelsoups}. We predict its transfer \emph{recovery} is
\textbf{monotone in the movement mediator} of \Cref{eq:movement}: composition helps \emph{most} for the
high-movement arms (SFT/ORPO --- much transfer lost, much to recover) and \emph{least} for the anchored
arms (DPO/KTO/GRPO --- little lost, little to recover, and averaging may even dilute their modest
represented gain). Doing this \emph{as a function of the training objective} is, to our knowledge,
unexplored.

\paragraph{Honest caveat on the mechanism.} The KL/on-policy anchoring account attributes low forgetting
partly to \emph{on-policy sampling keeping updates near the base}, not solely to an explicit KL penalty
term; and a \emph{static} preference-stage KL anchor can fail to prevent forgetting under stage-to-stage
data shift. So the anchored-static arms (DPO/KTO) and the on-policy arm (GRPO) may behave differently
\emph{even at matched movement}. We therefore report $M$ as a \emph{predictor}, not a proven cause, and
interpret the static-anchored and on-policy arms separately.

\subsection{Analysis plan and decision rules}
\label{sec:actII-analysis}

The analysis reuses the Act I / Act III machinery \emph{unchanged}, which is what makes the arms
comparable and the effort mostly eval-time:
\begin{itemize}[leftmargin=1.4em,itemsep=1pt]
  \item \textbf{Primary estimand.} Tie-aware macro average precision (the Act I definition) on
  represented-source and on dataset-held-out transfer, per (checkpoint, objective, seed); the paired
  base$\to$objective delta; and the fixed-panel aggregate. Reported as observed points with two-sided
  \emph{95\% paired hierarchical-bootstrap intervals} (10{,}000 replicates; family $+$ seed resampling),
  keeping each adapter paired with its own base. Descriptive and conditional on the fixed panel --- no
  significance test.
  \item \textbf{Mediator.} The base$\leftrightarrow$tuned movement of \Cref{eq:movement}, reported per
  (base, objective, seed), so H1/H2 can be plotted against it.
  \item \textbf{Composition (H2).} The existing base$+$adapter operator (\Cref{eq:comp}) is scored at
  analysis time for \emph{each} objective (nearly free, no extra training); recovery is
  (composition~$-$~objective) transfer AP regressed on the movement mediator.
  \item \textbf{Heterogeneity.} Per-checkpoint results reported; no post-hoc subgroup is promoted to a
  headline.
\end{itemize}

\paragraph{What would falsify each claim.} \textbf{H1} is falsified if the transfer ordering does not
track base-movement (e.g.\ the least-moving objective loses the most transfer), or if the arms are
statistically indistinguishable on transfer. The \textbf{GRPO null} is falsified --- a \emph{surprising
positive} we would report as such, not fold away --- if GRPO's transfer AP exceeds SFT's by an interval
strictly above zero on the fixed-panel aggregate. \textbf{H2} is falsified if composition recovery is
unrelated to, or inversely related to, base-movement across objectives.

\begin{takeaway}{What Act II will and will not establish}
Even fully run, this is a \emph{retrospective, estimation-only} description of one fixed panel with a
manifest inspected during Paper A development --- a locked rerun cannot make inspected data prospective, so
it is \emph{not} confirmatory. Nor is it a general ``RL vs.\ SFT'' verdict: the GRPO null is bounded to a
single-token binary guard on four compact checkpoints and one manifest, and is \emph{consistent with}, not
a refutation of, multi-token RL-generalization results. Intervals are conditional on the panel; no causal or
universal claim is licensed.
\end{takeaway}

\subsection{Results (pending the locked run)}
\label{sec:actII-results}

\begin{pending}{\textsc{pending} --- results table to be \code{\textbackslash input} as \code{tab:paper-c}}
No numbers are reported here. Once the locked Paper C run completes and its text-free per-row scores are
committed, the per-(checkpoint, objective) macro-AP table --- represented and transfer, base$\to$objective
deltas with 95\% hierarchical-bootstrap intervals, the fixed-panel aggregate, and the movement mediator
--- together with the objective$\times$composition (H2) recovery, will be generated from those scores and
inserted at the label \code{tab:paper-c}, replacing this box; the prose above stands unchanged as the
pre-registered design.
\end{pending}
% % PENDING: replaced by the real objective-axis table once the locked Paper C run completes.
\begin{center}
\fcolorbox{black!45}{black!4}{\begin{minipage}{0.92\linewidth}\small
\textbf{Results pending the locked run.} The objective-axis results table (SFT vs.\ DPO vs.\ GRPO
macro-AP on represented-source and dataset-held-out transfer, with per-checkpoint deltas and the
base$\leftrightarrow$tuned movement mediator) will be \texttt{\textbackslash input} here once the
pre-registered run in \Cref{sec:actII-analysis} completes and its per-row scores are committed. No
numbers are shown until then.
\end{minipage}}
\end{center}
  % <- provides \label{tab:paper-c}; add once the locked Paper C run completes.
